\documentclass[runningheads]{llncs}

\usepackage[T1]{fontenc}
\usepackage{array}
\usepackage{booktabs}
\usepackage{graphicx}
\usepackage{longtable}
\usepackage{pgfplots}
\usepackage{pgfplotstable}
\usepackage{url}

\newcommand{\twolinecol}[2]{\texttt{#1}\newline\texttt{#2}}
\pgfplotsset{compat=1.18}

\title{KNIME Workflows in Scientific Studies:
Publication Practices and Reproducibility Risks}
\titlerunning{KNIME Workflow Publication Practices}

\author{Vitalii Kaplan\orcidID{0009-0009-8181-2863}}
\authorrunning{V. Kaplan}
\institute{\email{2333275007@alu.istinye.edu.tr}}

\begin{document}
\maketitle

\begin{abstract}
In this paper, we study how visual workflows are reported in scientific
publications and how incomplete workflow preservation can make older studies
harder to reproduce as workflow platforms evolve. We use KNIME as a case
study, combining OpenAlex bibliometrics, a 100-record top-cited article audit
registry, workflow retrieval and current-KNIME opening/execution checks, and
repository-level mining of KNIME node metadata. OpenAlex shows that KNIME is
visible across a broad scientific literature, while the article assessment
shows that influential studies often report KNIME through figures,
descriptions, or tool mentions rather than downloadable workflow artifacts. In
the 100-record article audit registry, twenty-eight records report
downloadable or linked KNIME workflow files. Recorded retrieval attempts
obtained workflow artifacts for twelve article records. Workflows from those
twelve records opened in a current local KNIME environment; four article
records had at least one workflow execute successfully.
Repository mining gives the source-evolution context: deprecated ordinary
nodes increased from 14.83\% of registered ordinary nodes in the 2018 source
baseline to 33.33\% in the June 28, 2026 local source snapshot. Together,
these results show that weak workflow preservation and evolving node metadata
can make long-term reproduction of KNIME-based studies harder over time.

\keywords{Workflow reproducibility \and KNIME \and Deprecated nodes \and
Software evolution}
\end{abstract}

\section{Introduction}

KNIME is a widely used open-source platform for constructing scientific
data-analysis workflows, with more than a decade of use across published
studies. Its visual workflows represent analyses as connected nodes with
stored settings, data ports, and extension dependencies, which can make
computational work easier to inspect, reuse, and repeat. This representation
provides more operational detail than a narrative method description, but it
does not remove the need to preserve the executable workflow artifact and its
context. Reproduction requires access to the workflow file, KNIME version,
extensions, data, and execution environment. Figures and textual summaries do
not provide sufficient substitutes for these materials.

This issue matters for visual workflow platforms in general, because their
reproducibility value depends on preserving executable workflow artifacts and
the software context needed to interpret them. We study KNIME because it is an
open-source visual workflow platform with a long public development history and
substantial visibility in scientific literature. Our OpenAlex query returned
963 article-type records whose titles or abstracts match \texttt{KNIME}. The
matching records span platform papers, extension papers, tool-comparison
papers, and domain studies that use KNIME as an analysis interface. KNIME is
therefore a suitable case for studying whether visual workflows are preserved
as executable artifacts or only reported through descriptions and figures.

Our top-cited article collection covers the 100 most-cited KNIME-matching
records, with not-assessed placeholders retained where local full text is
unavailable. The structured audit reported in
Table~\ref{tab:top-cited-assessment} shows that publication practice is the
weak link. Twenty-eight records report downloadable or linked KNIME workflow
files, but workflow artifacts were obtained for only twelve article records in
the recorded retrieval attempts; other links were stale, blocked, incomplete,
or did not lead to a confirmed KNIME workflow artifact. Workflows from all
twelve obtained article records opened in the tested local KNIME environment,
and four article records had at least one workflow execute successfully. For
other influential studies, the strongest workflow evidence is usually weaker
than an executable artifact. Several records present screenshots or figures of
connected KNIME nodes, or describe workflow steps in text, without preserving
all executable workflow context.

These reporting practices are risky because KNIME continues to evolve. A
workflow that was executable when a paper was written may later depend on nodes
that are deprecated, hidden, migrated, removed, or supplied by third-party
extensions that are no longer easy to install. This paper therefore connects
publication evidence with platform-evolution evidence. We combine OpenAlex
bibliometrics, top-cited article retrieval and assessment, workflow-artifact
retrieval and current-version execution attempts, and repository-level mining
of KNIME node metadata. The goal is to show why publication practice and
source-code evolution must be considered together when evaluating long-term
reproducibility of KNIME-based studies.

The rest of the paper is organized as follows. Section~2 reviews related work
on computational reproducibility, open-source workflow systems, environment
preservation, and platform evolution. Section~3 presents the methods before
the data because the study depends on how the bibliometric records, article
assessments, and repository snapshots were collected. Section~4 describes the
resulting datasets. Section~5 reports the bibliometric, article-audit,
workflow-retrieval, manual compatibility-test, and repository-mining results.
Section~6 discusses
the implications, limitations, and future extensions of the study.

\section{Literature Review}

Reproducibility is a central requirement for computational research because
published results increasingly depend on software, data transformations, and
execution environments. Provenance research makes this requirement explicit:
visual analytics systems need records of data, analytical processes, and human
interpretation to make later inspection and reuse possible
\cite{Walker2013ProveML}. Studies of executable research artifacts show the
same problem in practice. Samuel and Mietchen found that biomedical Jupyter
notebooks often fail during dependency installation or re-execution even when
the notebooks themselves are available
\cite{SamuelMietchen2024JupyterReproducibility}.

Open-source scientific software helps reproducibility by making algorithms,
interfaces, and execution infrastructure inspectable. Weka, for example, is
presented as an extensible open-source data-mining environment with graphical
interfaces, reusable configurations, and plugin mechanisms
\cite{Hall2009WEKA}. KNIME is similarly used as an open-source visual workflow
platform, and recent GIScience work explicitly proposes KNIME-based workflows
as an approach for advancing replicable and reproducible research
\cite{Fu2025GIScienceKNIME}. However, openness of the platform alone is
insufficient. A workflow-based study also requires the executable workflow
file, input data, platform version, installed extensions, and other software
dependencies. Container research makes the same point at the environment
level: scientific computations are easier to reproduce when complete software
environments can be preserved and moved between systems
\cite{Kurtzer2017Singularity}. More broadly, the testing literature shows that
executable artifacts can remain sensitive to operating systems, dependency
versions, timing, external services, and other environmental factors
\cite{Rasheed2023FlakinessReview}.

Even when an executable workflow is published, long-term reproduction can be
affected by the evolution of the workflow platform itself. KNIME's node extension
schema states that a deprecated node is no longer shown in the node repository,
but is still loaded in existing workflows
\cite{KNIMEWorkbenchNodeSchema2026}. This design supports backward
compatibility, while also marking a change in the ordinary authoring surface.
KNIME workflows also persist node factory class names, and the platform
provides extension points for mapping old factory classes to newer
implementations and for registering workflow-node migration rules
\cite{KNIMECoreClassMapperSchema2026,KNIMEMigrationRuleSchema2026}. These
mechanisms show that continued workflow loading may depend on compatibility
metadata, extension availability, and migration support. We therefore treat
KNIME node deprecation as a measurable risk marker for long-term workflow
reproducibility.

\section{Methods}

All scripts and data used for this study, except the other articles assessed
as study objects, are available in the public replication repository
\cite{KaplanReproducibilityRisksArticle2026}. Relative paths reported in this
paper should be interpreted as paths within that GitHub repository.

\subsection{Bibliometric Evidence}

We collected bibliometric evidence from the OpenAlex Works API
\cite{Priem2022OpenAlex} on June 29, 2026. The query used OpenAlex
title-and-abstract search for the word \texttt{KNIME} and restricted the
result set to OpenAlex records whose type is \texttt{article}. The same query
can be reproduced in the
OpenAlex web interface by selecting ``Search title and abstract'', entering
\texttt{KNIME}, and filtering the result type to \texttt{article}. The
corresponding web-interface request can be represented as:

\begin{quote}
\path{https://openalex.org/works}\\
\path{page=1}\\
\path{sort=relevance_score:desc}\\
\path{search.title_and_abstract=KNIME}\\
\path{filter=type:article}
\end{quote}

The first API request, which returns the first 200 records, is reproducible as:

\begin{quote}
\path{https://api.openalex.org/works}\\
\texttt{?search.title\_and\_abstract=KNIME}\\
\texttt{\&filter=type:article}\\
\texttt{\&per-page=200}\\
\texttt{\&cursor=*}
\end{quote}

OpenAlex uses cursor pagination for large result sets. After the first request,
the collector reads \texttt{meta.next\_cursor} from each response and submits
the same request with \texttt{cursor=<next-cursor>} until no further results
are returned. The value \texttt{per-page=200} is the maximum page size used
here.

The collection script is
\path{scripts/collect_openalex_knime_works.py}.
It stores raw and processed collection artifacts as follows:
\begin{itemize}
\item raw OpenAlex response pages:
\path{data/original/openalex/pages/}
\item complete record stream:
\path{data/original/openalex/works.jsonl}
\item compact CSV view:
\path{data/original/openalex/works.csv}
\item endpoint, query parameters, collection time, result count, and page
manifest:
\path{data/original/openalex/metadata.json}
\end{itemize}
The current run collected 963 records across 5 OpenAlex response pages.
Derived bibliometric tables are generated from the JSONL file by
\path{scripts/build_openalex_knime_bibliometrics.py}.

\subsection{Top-Cited Article Assessment}

For the article-level audit, we sorted the processed OpenAlex article records
by \texttt{cited\_by\_count} in descending order and selected the 100
most-cited records. Open-access papers were downloaded directly. Papers that
required access through a publisher were downloaded where institutional
subscription access was available or were added manually when a local copy was
available. Not all 100 papers could be downloaded; records without local full
text were retained as not-assessed placeholders so that the denominator stayed
visible.

Downloaded PDFs were converted to plain text with
\path{scripts/extract_article_texts.py}. Where extracted text followed a
two-column reading order, we normalized it to a one-column reading copy with
\path{scripts/normalize_article_text_columns.py}; this step used script output
and LLM-assisted correction checks. We then processed the article text into a
structured audit stored in
\path{data/processed/audit/knime_most_cited_article_assessments.json}. Each
record contains descriptive metadata, full-text access status, the article's
relation to KNIME, and reproducibility-relevant flags for version reporting,
workflow artifacts, workflow figures or text descriptions, input data, code,
and extension information.

The audit fields are defined in
\path{data/processed/audit/knime_article_audit_questions.json}. Positive
article-text flags require direct quote support from the processed text,
including line references, article section, and an audit note. The support and
provenance fields were refreshed with
\path{scripts/refresh_article_audit_support.py}, and article tables were
generated from the JSON audit with
\path{scripts/build_article_audit_tables.py} and
\path{scripts/build_knime_use_workflow_reporting_table.py}. For records that
reported downloadable or linked KNIME workflows, we attempted to retrieve the
workflow artifacts, recorded the download outcomes in the workflow inventory,
and finally tested whether the downloaded workflows could be opened and
executed in a current local KNIME installation.

\subsection{Repository-Level Analysis}

We study KNIME node deprecation through repository mining of the public
\texttt{knime-oss} GitHub organization. KNIME Analytics Platform is a
workflow-based data-analysis system in which analyses are assembled from
connected processing nodes
\cite{Berthold2009KNIME,KNIMEUserGuide2026}. KNIME extensions are distributed
through update sites and extend the platform with additional functionality
\cite{KNIMEExtensions2026}. The KNIME source repositories were cloned locally,
and each source snapshot was reconstructed by checking out every non-hidden
top-level Git repository to the latest commit at or before a selected target
date. This checkout step is implemented by
\path{scripts/checkout_knime_oss_by_date.sh}.
The script records a manifest for each snapshot, including repository name,
checkout status, target date, selected commit, selected commit date, and
previous head.

The current repository-mining corpus contains 91 immediate Git repositories in
the local clone. The date-based mining process includes 90 non-hidden immediate
repositories. The hidden \texttt{.github} repository is excluded because it
contains organization metadata rather than KNIME node implementation code. The
exact repository list is stored in
\path{data/processed/knime_snapshots/}
\path{knime_oss_repositories.csv}.

For each date-based source snapshot, we parse KNIME metadata files rather than
grepping source text. The main deprecation marker is
\texttt{deprecated="true"} on ordinary node registrations in
\texttt{plugin.xml} under the Eclipse extension point
\begin{quote}
\texttt{org.knime.workbench.repository.nodes}
\end{quote}
The KNIME schema states that deprecated nodes are no longer shown in the node
repository but are still loaded in existing workflows
\cite{KNIMEWorkbenchNodeSchema2026}. Dynamic node sets use a separate
extension point:
\begin{quote}
\texttt{org.knime.workbench.repository.nodesets}
\end{quote}
We count these registrations separately. Node-description XML files ending in
\texttt{NodeFactory.xml} and rooted at \texttt{knimeNode} are parsed as a
secondary deprecation source.
Hidden nodes, marked with \texttt{hidden="true"}, are counted separately from
deprecated nodes.

The extraction also records compatibility-related metadata from two extension
points:
\begin{quote}
\texttt{org.knime.core.NodeFactoryClassMapper}
\end{quote}
\begin{quote}
\texttt{org.knime.workflow.migration.NodeMigrationRule}
\end{quote}
Together, these extension points document mechanisms for mapping persisted
node-factory class names to newer implementations and for registering
workflow-node migration rules
\cite{KNIMECoreClassMapperSchema2026,KNIMEMigrationRuleSchema2026}. These
records are not counted as deprecation markers, but they are relevant for later
analysis of whether deprecated or replaced nodes have explicit migration
support.

The parser is implemented in
\path{scripts/collect_knime_node_snapshot.py}.
It uses Python's standard XML parser and skips \texttt{.git},
\texttt{target}, \texttt{bin}, and \texttt{.metadata} directories. Boolean XML
attributes are treated as true only when their value is case-insensitively
equal to \texttt{true}.

For each snapshot, the extractor writes per-record tables for node
registrations, node descriptions, factory-class mappers, migration rules, and a
local summary. The per-snapshot outputs are stored under
\path{data/original/knime_snapshots/<snapshot-date>/}
and include:
\begin{itemize}
\item \path{logs/checkout_<snapshot-date>.csv}
\item \path{plugin_nodes.csv}
\item \path{node_descriptions.csv}
\item \path{factory_class_mappers.csv}
\item \path{migration_rules.csv}
\item \path{summary.csv}
\end{itemize}
The cross-snapshot table is stored as:
\begin{quote}
\path{data/processed/knime_snapshots/}

\path{knime_node_snapshot_summary.csv}
\end{quote}
It is generated from these per-snapshot outputs by
\path{scripts/build_knime_node_snapshot_summary.py}.

The summary builder aggregates ordinary node registrations, dynamic nodesets,
unique factory classes, deprecated nodes, hidden nodes, node-description files,
deprecated node descriptions, factory-class mappers, and migration rules. It
also counts processed and skipped repositories from each checkout manifest.
Transition columns compare adjacent chronological snapshots. Node identity is
based on \texttt{factory\_class} where available. Otherwise, the fallback key
is \texttt{plugin\_xml:element:category\_path}. These transition counts are
therefore metadata-level approximations and are strongest for nodes with stable
factory classes.

All project Python scripts are intended to be run with Python 3.11 or newer.
The current scripts use only the Python standard library. No third-party
Python packages are required for the OpenAlex and KNIME repository-mining
pipelines.

\section{Data}

\subsection{Bibliometric Evidence}

The OpenAlex bibliometric data are organized into original API responses and
processed analysis tables. The original data are stored under
\path{data/original/openalex/}.
They include raw response pages, a complete JSONL export, a compact CSV export,
and query metadata. These files preserve the article-level records used to
measure the size, time trend, venues, fields, and accessibility of the
KNIME-matching literature.

Derived bibliometric tables are stored under
\path{data/processed/openalex/}.
Their roles in the analysis are:

\begin{itemize}
\item \path{openalex/openalex_knime_articles.csv}: compact article-level
metadata used as the main processed bibliography table.
\item \path{openalex/openalex_knime_articles_by_year.csv}: annual publication
counts used to describe the time trend of KNIME-matching literature.
\item \path{openalex/openalex_knime_top_venues.csv}: source-level aggregation
used to identify journals, proceedings, repositories, and platforms where the
matching records appear.
\item \path{openalex/openalex_knime_top_domains.csv}: OpenAlex domain
aggregation used to describe the broad disciplinary distribution.
\item \path{openalex/openalex_knime_top_fields.csv}: OpenAlex field
aggregation used to describe the finer disciplinary distribution.
\item \path{openalex/openalex_knime_most_cited.csv}: citation-ranked subset
used to select influential records for later manual assessment.
\item \path{openalex/openalex_knime_likely_workflow_platform.csv}: heuristic
subset used to prioritize papers likely to use KNIME as a workflow platform.
\end{itemize}

The processed OpenAlex directory also contains additional summary files:
\begin{itemize}
\item \path{openalex/openalex_knime_open_access_availability.json}
\item \path{openalex/openalex_knime_role_classification_summary.json}
\end{itemize}
These files summarize artifact-access signals and role-classification counts.

\subsection{Top-Cited Article Assessment}

The top-cited article dataset records the structured assessment of the 100
most-cited OpenAlex records. The public replication package contains processed
assessment metadata and extracted text, but not the article PDFs. The
citation-ranked source list is
\path{data/processed/openalex/openalex_knime_most_cited.csv}. Download logs
and text-conversion manifests are stored under
\path{data/processed/audit/logs/}. Processed one-column reading files are
stored under \path{data/processed/articles/}.

The main assessment file is
\path{data/processed/audit/knime_most_cited_article_assessments.json}. Each
record stores the OpenAlex rank, article metadata, full-text access status,
KNIME role, reported KNIME version, workflow-artifact status, data and code
availability, linked resources visible in the paper, quote-backed support for
positive article-text flags, and short evidence notes. Records without
accessible full text are retained as not-assessed placeholders.

The article-level summary tables are generated from the assessment JSON rather
than edited manually. Table~\ref{tab:top-cited-assessment} is generated from
the audit flags, and Table~\ref{tab:knime-use-assessment} combines the
same assessment file with the workflow-retrieval inventory. Downloadable
workflow references, download outcomes, local workflow paths, and manual
opening or execution notes are recorded in
\path{data/processed/audit/knime_downloadable_workflow_references.json}.
Manual KNIME-opening screenshots are stored under each workflow directory's
\path{opened/} subdirectory.

\subsection{Repository-Level Analysis}

The KNIME repository-mining data are organized as original per-snapshot audit
files and one processed cross-snapshot summary. The processed summary is stored
as:
\begin{quote}
\path{data/processed/knime_snapshots/}

\path{knime_node_snapshot_summary.csv}.
\end{quote}
Per-snapshot audit files are stored under
\path{data/original/knime_snapshots/<snapshot-date>/} and include:

\begin{itemize}
\item \path{logs/checkout_<snapshot-date>.csv}
\item \path{plugin_nodes.csv}
\item \path{node_descriptions.csv}
\item \path{factory_class_mappers.csv}
\item \path{migration_rules.csv}
\item \path{summary.csv}
\end{itemize}

\begingroup
\small
\begin{longtable}{%
  >{\raggedright\arraybackslash}p{0.40\textwidth}
  >{\raggedright\arraybackslash}p{0.18\textwidth}
  >{\raggedright\arraybackslash}p{0.32\textwidth}}
\caption{Columns in the processed KNIME node snapshot summary.}
\label{tab:summary-columns}\\
\toprule
Name & Type & Short description \\
\midrule
\endfirsthead
\toprule
Name & Type & Short description \\
\midrule
\endhead
\bottomrule
\endfoot
\texttt{snapshot\_id} & string &
Stable identifier for the snapshot. \\
\texttt{snapshot\_kind} & categorical &
Snapshot basis, currently date-based. \\
\texttt{snapshot\_date} & date &
Target date used for repository checkout. \\
\texttt{knime\_version} & string &
Version label or source-date context assigned by the summary script. \\
\texttt{source\_basis} & categorical &
Method used to obtain the source snapshot. \\
\texttt{repos\_processed} & integer &
Repositories with a checkout commit in the snapshot manifest. \\
\texttt{repos\_skipped} & integer &
Repositories without a commit at or before the target date. \\
\texttt{plugin\_xml\_files} & integer &
Distinct parsed \texttt{plugin.xml} files contributing node metadata. \\
\texttt{repos\_with\_plugin\_xml} & integer &
Repositories contributing at least one parsed \texttt{plugin.xml}. \\
\texttt{registered\_nodes} & integer &
Ordinary node registrations in the node extension point. \\
\texttt{registered\_nodesets} & integer &
Dynamic nodeset registrations in the nodeset extension point. \\
\texttt{registered\_total} & integer &
Sum of ordinary node and nodeset registrations. \\
\texttt{unique\_factory\_classes} & integer &
Distinct non-empty factory classes among ordinary node registrations. \\
\texttt{deprecated\_nodes} & integer &
Ordinary nodes with \texttt{deprecated="true"}. \\
\texttt{deprecated\_nodesets} & integer &
Nodesets with \texttt{deprecated="true"}. \\
\texttt{deprecated\_total} & integer &
Sum of deprecated ordinary nodes and deprecated nodesets. \\
\twolinecol{unique\_deprecated}{\_factory\_classes} & integer &
Distinct factory classes among deprecated ordinary nodes. \\
\texttt{deprecated\_node\_percent} & percent &
Deprecated ordinary nodes divided by registered ordinary nodes. \\
\texttt{hidden\_nodes} & integer &
Ordinary nodes with \texttt{hidden="true"}. \\
\texttt{hidden\_node\_percent} & percent &
Hidden ordinary nodes divided by registered ordinary nodes. \\
\twolinecol{deprecated\_and}{\_hidden\_nodes} & integer &
Ordinary nodes marked both deprecated and hidden. \\
\texttt{node\_description\_files} & integer &
Parsed \texttt{*NodeFactory.xml} files rooted at \texttt{knimeNode}. \\
\twolinecol{description\_deprecated}{\_files} & integer &
Node-description files with \texttt{deprecated="true"}. \\
\twolinecol{description\_deprecated}{\_percent} & percent &
Deprecated node descriptions divided by parsed node descriptions. \\
\twolinecol{factory\_class\_mapper}{\_count} & integer &
Registered \texttt{NodeFactoryClassMapper} contributions. \\
\texttt{migration\_rule\_count} & integer &
Registered \texttt{NodeMigrationRule} contributions. \\
\twolinecol{nodes\_added}{\_since\_previous} & integer/blank &
Node identity keys present in this snapshot but not the previous one. \\
\twolinecol{nodes\_removed}{\_since\_previous} & integer/blank &
Node identity keys present in the previous snapshot but absent here. \\
\twolinecol{nodes\_newly\_deprecated}{\_since\_previous} & integer/blank &
Common node keys that became deprecated since the previous snapshot. \\
\twolinecol{nodes\_no\_longer\_deprecated}{\_since\_previous} &
integer/blank &
Common node keys no longer marked deprecated. \\
\twolinecol{nodes\_newly\_hidden}{\_since\_previous} & integer/blank &
Common node keys that became hidden since the previous snapshot. \\
\twolinecol{nodes\_no\_longer\_hidden}{\_since\_previous} & integer/blank &
Common node keys no longer marked hidden. \\
\twolinecol{nodes\_category\_changed}{\_since\_previous} & integer/blank &
Common node keys whose observed category-path set changed. \\
\end{longtable}
\endgroup

\section{Results}

\subsection{Bibliometric Evidence}

The OpenAlex title-and-abstract query returned 963 article-type records. This
is the first piece of evidence for the study: KNIME is visible in a sizeable
scientific literature, so preservation of KNIME workflows is not only a niche
platform-maintenance issue. The annual series indicates that KNIME-matching
literature becomes visible in OpenAlex from the mid-2000s and grows into a
sustained publication stream after 2015 (Fig.~\ref{fig:openalex-by-year}). The
highest annual count in the current export is 104 records in 2023. Counts for
2026 are incomplete because the query was collected on June 29, 2026.

\begin{figure}
\centering
\begin{tikzpicture}
\begin{axis}[
  width=0.95\textwidth,
  height=0.34\textwidth,
  ybar,
  bar width=3.5pt,
  xlabel={Publication year},
  ylabel={Article records},
  ymin=0,
  xmin=2005,
  xmax=2027,
  enlarge x limits=0.02,
  xtick={2006,2010,2015,2020,2026},
  ymajorgrids=true,
  grid style={gray!25},
  tick label style={font=\scriptsize},
  label style={font=\small},
]
\addplot table[
  x index=0,
  y index=1,
  col sep=comma,
  ignore chars={\"}
] {../data/processed/openalex/openalex_knime_articles_by_year.csv};
\end{axis}
\end{tikzpicture}
\caption{OpenAlex article records matching \texttt{KNIME} in title or
abstract by publication year.}
\label{fig:openalex-by-year}
\end{figure}

The domain distribution shows that the KNIME-matching literature is not
confined to one disciplinary area (Table~\ref{tab:openalex-domains}). OpenAlex
assigns just over half of the records to Physical Sciences, but the remaining
records are spread across Social Sciences, Life Sciences, and Health Sciences.
This matters for the study because workflow reproducibility risk is not only a
software-engineering concern. If KNIME workflows are used across computational
chemistry, biomedical analysis, education, business analytics, and related
areas, then missing workflow artifacts and changing node implementations can
affect published computational evidence in several research communities.

\begin{table}
\centering
\caption{OpenAlex domains for KNIME title-and-abstract matches.}
\label{tab:openalex-domains}
\begin{tabular}{lrr}
\toprule
Domain & Articles & Share \\
\midrule
Physical Sciences & 505 & 52.44\% \\
Social Sciences & 187 & 19.42\% \\
Life Sciences & 149 & 15.47\% \\
Health Sciences & 122 & 12.67\% \\
\bottomrule
\end{tabular}
\end{table}

The ten largest OpenAlex fields are led by Computer Science (350),
Biochemistry, Genetics and Molecular Biology (115), Medicine (92),
Engineering (71), and Social Sciences (66). These field counts show why a
paper about KNIME reproducibility cannot focus only on the KNIME platform
literature. The central reproducibility question is how scientific studies
that use KNIME as a tool report the executable workflow, data, versions, and
extensions needed for later reproduction.

Open-access metadata indicate 619 open-access records, 625 records with some
full-text availability, and 421 records with a PDF signal. These signals are
important because later workflow audits require access to paper text,
supplements, repositories, and sometimes PDF-only workflow descriptions. The
counts are only availability indicators. They do not by themselves establish
that a workflow artifact, KNIME version, input data, or extension list is
available.

The heuristic role classification separates likely papers about KNIME or KNIME
extensions from papers that likely use KNIME as an analysis tool. The current
rules classify 156 records as about KNIME or a KNIME extension, 314 as using
KNIME as a tool, 28 as about or mentioning KNIME in the title, and 465 as
ambiguous OpenAlex matches. The same heuristic marks 410 records as likely
using KNIME as a workflow platform. These labels are useful for triage: they
identify a substantial candidate set where workflow reporting can matter, but
final case-study claims require manual full-text and artifact assessment.

\subsection{Top-Cited Article Assessment}

We built a top-cited retrieval pool from OpenAlex records matching
\texttt{KNIME}. The statistics in this subsection use the 100-record structured
audit registry. The purpose of this part of the study is to ask whether
KNIME-matching papers preserve the evidence needed to open and reproduce a
KNIME-based workflow: full text, KNIME-use role, version information, workflow
files, workflow figures or textual descriptions, data, code, and extension
context.

The audit separates background and platform papers about KNIME itself from
workflow-reporting KNIME use cases. Eight of the 100 records are about KNIME
itself or a KNIME extension. These records are important for understanding the
KNIME ecosystem, but they do not need the same artifact-preservation audit as
domain studies that use KNIME. We therefore retain their descriptive metadata
but leave their workflow-reporting flag maps empty.
Table~\ref{tab:top-cited-assessment} uses all 100 records in the
registry as the denominator. It first reports the classification of the
records, then reports the workflow-reporting flags observed in the structured
audit.

Table~\ref{tab:top-cited-assessment} summarizes the statistics-ready flags in
the structured audit file. Sixty-seven records use KNIME as a workflow, tool,
interface, or implementation context. Fifty-nine records contain textual
workflow or node descriptions, and thirty-seven show workflow screenshots or
figures. Twenty-eight records report downloadable or linked KNIME workflow
files, and the recorded retrieval attempts obtained workflow artifacts for
twelve of these article records. Version and dependency reporting remain limited: eighteen
records report a KNIME version, thirty report KNIME extension or plugin
dependencies, and nineteen report an installation source. Data and code signals
are more common than retrievable workflow-artifact evidence: fifty-five records
report input-data availability, twenty-three provide a direct input-data
resource, and seventeen provide code or scripts. In this study, we do not
verify data availability itself; we only record whether the article text
reports it. The
same limitation applies to code or script availability.
This pattern is important because a reader may find data or code without being
able to recover the actual KNIME workflow, node settings, KNIME version, or
extension context.

\begingroup
\scriptsize
\renewcommand{\arraystretch}{1.18}
\begin{longtable}{|%
  >{\raggedright\arraybackslash}p{0.46\textwidth}|
  r|
  r|
  >{\raggedright\arraybackslash}p{0.25\textwidth}|}
\caption{Flag summary for the top-cited article audit. Shares use all
100 records in the registry as the denominator.}
\label{tab:top-cited-assessment}\\
\toprule
Audit category or flag & Count & Share of 100 & Interpretation \\
\midrule
\endfirsthead
\toprule
Audit category or flag & Count & Share of 100 & Interpretation \\
\midrule
\endhead
\bottomrule
\endfoot
Audit records & 100 & 100.0\% & Most-cited KNIME-matching article records selected from OpenAlex \\
Assessed from full text & 82 & 82.0\% & Article text available for manual assessment \\
Not assessed from full text & 18 & 18.0\% & Full text unavailable at assessment time \\
About KNIME & 8 & 8.0\% & Background/platform or extension papers, not deeper workflow-reporting cases \\
Not a KNIME use case & 7 & 7.0\% & KNIME appears in the record but is not a KNIME-based study \\
\midrule
Uses KNIME & 67 & 67.0\% & KNIME used as a workflow, tool, interface, or implementation context \\
Workflow or nodes described in text & 59 & 59.0\% & Named workflow steps, nodes, modules, or components recorded \\
Workflow screenshots or figures & 37 & 37.0\% & Workflow shown visually in article or supplement \\
Reports KNIME version & 18 & 18.0\% & Specific KNIME version value found \\
Reports downloadable workflows & 28 & 28.0\% & Executable workflow files provided or linked \\
Reports extension/plugin dependencies & 30 & 30.0\% & KNIME extension or node-library dependency reported \\
Reports extension installation source & 19 & 19.0\% & Update site, project URL, or installation source reported \\
\midrule
Provides code or scripts & 17 & 17.0\% & Code repository, scripts, or software code reported \\
Reports input-data availability & 55 & 55.0\% & Data availability stated, with or without direct URL \\
Direct input-data resource & 23 & 23.0\% & Direct data URL or resource pointer recorded \\
\end{longtable}
\endgroup

\begingroup
\scriptsize
\renewcommand{\arraystretch}{1.18}
\begin{longtable}{|%
  >{\raggedright\arraybackslash}p{0.62\textwidth}|
  r|
  r|}
\caption{Workflow-reporting and retrieval signals among the 67 assessed KNIME-use cases.}
\label{tab:knime-use-assessment}\\
\toprule
Audit flag & Count & Share of KNIME-use cases \\
\midrule
\endfirsthead
\toprule
Audit flag & Count & Share of KNIME-use cases \\
\midrule
\endhead
\bottomrule
\endfoot
Uses KNIME & 67 & 100.0\% \\
Workflow or nodes described in text & 59 & 88.1\% \\
Workflow screenshots or figures & 37 & 55.2\% \\
Reports KNIME version & 18 & 26.9\% \\
Reports extension/plugin dependencies & 30 & 44.8\% \\
Reports extension installation source & 19 & 28.4\% \\
Reports downloadable workflows & 28 & 41.8\% \\
Articles with successfully downloaded workflows & 12 & 17.9\% \\
\end{longtable}
\endgroup

Among the 28 article records that report downloadable or linked KNIME workflow
files, Table~\ref{tab:workflow-retrieval-results} summarizes the 19 records
for which we recorded an artifact-retrieval outcome. In this table, the
\emph{Success} column marks whether a KNIME workflow artifact was actually
obtained: \textbf{X} means that a workflow file, archive, or workflow
directory was retrieved, while \textbf{-} means that no KNIME workflow artifact
was obtained or confirmed. Workflow artifacts or workflow directories were
obtained for 12 article records. The remaining recorded retrieval attempts
failed because links were obsolete, blocked, unresolved, or led to repositories
or supplements that did not contain a confirmed KNIME workflow file.

Table~\ref{tab:workflow-opening-results} then reports the manual current-KNIME
opening and execution checks for those 12 downloaded article records. Here, the
\emph{Success} column has a stricter meaning: \textbf{X} means that at least
one workflow for the article opened and executed successfully in the tested
environment, while \textbf{-} means that the workflow opened but did not have a
confirmed successful end-to-end run, or failed during execution. Workflows from
all 12 obtained article records opened in the local KNIME environment (KNIME
5.8.3 LTS). Four article records had at least one workflow execute
successfully: PAINS, Webinar Pricing Analytics, ImageJ ecosystem integration,
and high-content organelle trafficking
\cite{Saubern2011PAINS,VillarroelOrdenes2021KNIMEHub,Dietz2020ImageJKNIME,Pal2018OrganelleKNIME}.
The remaining opened article records did not execute successfully in this
manual check. This does not show that those workflows are unreproducible; they
may require additional configuration, dependencies, or repair before a
successful run. Across the opened workflows, the screenshots also show
deprecated or legacy nodes, but we treat this as qualitative compatibility
evidence here rather than counting visible deprecated nodes. The screenshots
are archived under each workflow directory's \texttt{opened/} subdirectory in
\path{data/original/workflows/}.

\begingroup
\scriptsize
\renewcommand{\arraystretch}{1.18}
\begin{longtable}{|%
  >{\centering\arraybackslash}p{0.04\textwidth}|
  >{\raggedright\arraybackslash}p{0.23\textwidth}|
  >{\centering\arraybackslash}p{0.09\textwidth}|
  >{\raggedright\arraybackslash}p{0.54\textwidth}|}
\caption{Workflow retrieval results for article records with attempted or
recorded artifact retrieval outcomes.}
\label{tab:workflow-retrieval-results}\\
\toprule
\# & Article record & Success & Retrieval note \\
\midrule
\endfirsthead
\toprule
\# & Article record & Success & Retrieval note \\
\midrule
\endhead
\bottomrule
\endfoot
1 &
PAINS workflows, DOI \path{10.1002/minf.201100076} &
\textbf{X} &
Automated myExperiment page and download requests returned HTTP 403 HTML, but
the user manually obtained both RDKit and Indigo workflow ZIP archives; both
contain \texttt{workflow.knime} files. \\
2 &
Chemical data curation workflow, DOI \path{10.1186/s13321-018-0315-6} &
\textbf{X} &
GitHub master archive was downloaded after retry and passed ZIP integrity
checking; KNIME workflow files were found in the archive. \\
3 &
ASD genes prediction workflow, DOI \path{10.1371/journal.pone.0208626} &
\textbf{X} &
GitHub archive passed ZIP integrity checking; an inner KNIME workflow archive
was extracted to a workflow directory. \\
4 &
Bitterness/sweetness classifier workflow, DOI \path{10.3389/fchem.2018.00093} &
\textbf{-} &
The original project page redirected to VirtualTaste, and the redirected page
returned HTTP 404; no workflow archive was retrieved. \\
5 &
Medicinal chemistry filters workflow, DOI \path{10.3390/encyclopedia3020035} &
\textbf{X} &
GitLab archive passed ZIP integrity checking and contained KNIME workflow
files; one KNIME Hub URL returned a page, while an older generic Hub URL
returned 404 JSON. \\
6 &
MS-ready structures repository, DOI \path{10.1186/s13321-018-0299-2} &
\textbf{-} &
The GitHub archive passed ZIP integrity checking, but the extracted repository
contained only \texttt{LICENSE} and \texttt{README.md}; no \texttt{.knwf} or
\texttt{workflow.knime} file was found. \\
7 &
KNIME-OMICS reference, DOI \path{10.1016/j.jprot.2016.08.002} &
\textbf{-} &
The referenced GitHub organization page returned HTTP 404; no workflow archive
was retrieved. \\
8 &
QSAR tools workflow, DOI \path{10.1080/07391102.2018.1456975} &
\textbf{-} &
The referenced \path{teqip.jdvu.ac.in/QSAR_Tools/} URL failed DNS resolution;
no workflow archive was retrieved. \\
9 &
Webinar Pricing Analytics workflow, DOI \path{10.1016/j.jbusres.2021.08.036} &
\textbf{X} &
The user manually downloaded the KNIME workflow file and added local KNIME
opening screenshots. \\
10 &
Verhaar scheme workflow, DOI \path{10.1016/j.chemosphere.2015.06.009} &
\textbf{-} &
The article reports a supplementary KNIME workflow; standard Elsevier
supplement URLs resolved as Excel files, but no direct KNIME workflow archive
URL was confirmed. \\
11 &
Spontaneous tail coiling workflow, DOI \path{10.1016/j.ntt.2020.106918} &
\textbf{X} &
The user manually downloaded the KNIME workflow file and added local KNIME
opening screenshots. \\
12 &
Caco-2 permeability prediction workflow, DOI \path{10.3390/pharmaceutics14101998} &
\textbf{X} &
The user manually downloaded the KNIME workflow file and added local KNIME
opening screenshots. \\
13 &
KniMet workflow, DOI \path{10.1007/s11306-018-1349-5} &
\textbf{X} &
The Zenodo DOI resolved to record 1196407; \texttt{KniMet-master.zip} was
downloaded and contains \texttt{KniMet.knwf}. \\
14 &
NGS sample workflow, DOI \path{10.1093/bioinformatics/btr478} &
\textbf{X} &
The user manually downloaded the myExperiment archive after automated requests
were blocked by HTTP 403; the ZIP contains a KNIME workflow directory with
saved data. \\
15 &
ImageJ ecosystem workflows, DOI \path{10.3389/fcomp.2020.00008} &
\textbf{X} &
Six KNIME Hub artifact metadata endpoints returned signed URLs and the
corresponding \texttt{.knwf} files were downloaded; one recorded reference
appears stale or unavailable. \\
16 &
Multi-template matching repository, DOI \path{10.1186/s12859-020-3363-7} &
\textbf{-} &
The GitHub archive was downloaded, but it contained documentation/images only
and no \texttt{.knwf}, \texttt{workflow.knime}, or KNIME workflow archive; a
guessed NodePit URL returned HTTP 404. \\
17 &
GediNET workflows, DOI \path{10.1038/s41598-022-24421-0} &
\textbf{X} &
GitHub master archive was downloaded and contains two \texttt{.knwf} files;
the recorded KNIME Hub short link could not be retrieved. \\
18 &
Nuclear receptor ligand workflow, DOI \path{10.1002/minf.201400188} &
\textbf{-} &
The article reports supplementary workflow material, but Wiley article and
guessed supplementary download endpoints returned Cloudflare challenge pages;
Crossref metadata did not expose a supplementary workflow URL. \\
19 &
High-content organelle trafficking workflow, DOI \path{10.1038/sdata.2018.241} &
\textbf{X} &
The figshare collection DOI resolved through the figshare API; the default
KNIME workflow file was downloaded as \texttt{KNIME workflow.zip}. \\
\end{longtable}
\endgroup

\begingroup
\scriptsize
\renewcommand{\arraystretch}{1.18}
\begin{longtable}{|%
  >{\centering\arraybackslash}p{0.04\textwidth}|
  >{\raggedright\arraybackslash}p{0.25\textwidth}|
  >{\centering\arraybackslash}p{0.09\textwidth}|
  >{\raggedright\arraybackslash}p{0.51\textwidth}|}
\caption{Manual opening and execution results for downloaded KNIME workflow
artifacts.}
\label{tab:workflow-opening-results}\\
\toprule
\# & Article record & Success & Manual KNIME result \\
\midrule
\endfirsthead
\toprule
\# & Article record & Success & Manual KNIME result \\
\midrule
\endhead
\bottomrule
\endfoot
1 &
PAINS workflows, DOI \path{10.1002/minf.201100076} &
\textbf{X} &
RDKit and Indigo workflow ZIPs were tested. At article-record level, at least
one workflow opened and executed successfully. The separate Indigo workflow
opened but was blocked by a missing Indigo KNIME Integration node. \\
2 &
Chemical data curation workflow, DOI \path{10.1186/s13321-018-0315-6} &
\textbf{-} &
The CompTox rebuilt workflow file opened, but execution failed inside a
metanode during the manual test. \\
3 &
ASD genes prediction workflow, DOI \path{10.1371/journal.pone.0208626} &
\textbf{-} &
The extracted ASD genes prediction workflow opened, but execution failed in an
R node because required R packages and functions were unavailable in the test
environment. \\
4 &
Medicinal chemistry filters workflow, DOI \path{10.3390/encyclopedia3020035} &
\textbf{-} &
The medicinal chemistry filters workflow file opened, but a successful
end-to-end execution was not established in the available notes. \\
5 &
Webinar Pricing Analytics workflow, DOI \path{10.1016/j.jbusres.2021.08.036} &
\textbf{X} &
The Webinar Pricing Analytics workflow file opened and executed successfully
in the current local KNIME environment. \\
6 &
Spontaneous tail coiling workflow, DOI \path{10.1016/j.ntt.2020.106918} &
\textbf{-} &
The spontaneous tail coiling measurement workflow file opened in the current
local KNIME environment, but a successful end-to-end execution was not
established in the available notes. \\
7 &
Caco-2 permeability prediction workflow, DOI \path{10.3390/pharmaceutics14101998} &
\textbf{-} &
The manual Caco-2 prediction \texttt{.knwf} file opened in the current local
KNIME environment, but a successful end-to-end execution was not established
in the available notes. \\
8 &
KniMet workflow, DOI \path{10.1007/s11306-018-1349-5} &
\textbf{-} &
The KniMet workflow from the Zenodo archive opened in the current local KNIME
environment, but successful execution was not established in the available
notes. \\
9 &
NGS sample workflow, DOI \path{10.1093/bioinformatics/btr478} &
\textbf{-} &
The NGS sample workflow archive with saved data opened in the current local
KNIME environment, but successful execution was not established in the
available notes. \\
10 &
ImageJ ecosystem workflows, DOI \path{10.3389/fcomp.2020.00008} &
\textbf{X} &
Six downloaded KNIME Hub \texttt{.knwf} files opened and executed normally in
the user environment after installing additional required KNIME/ImageJ-related
plugins. \\
11 &
GediNET workflows, DOI \path{10.1038/s41598-022-24421-0} &
\textbf{-} &
The GediNET GitHub archive opened in the current local KNIME environment, but
successful execution was not established in the available notes. \\
12 &
High-content organelle trafficking workflow, DOI \path{10.1038/sdata.2018.241} &
\textbf{X} &
\texttt{KNIME workflow.zip} opened and executed normally in the user
environment after installing additional required plugins. \\
\end{longtable}
\endgroup

\subsection{Repository-Level Analysis}

The repository-mining results explain why weak workflow reporting becomes more
costly over time. Deprecated nodes are not an isolated metadata corner case in
the KNIME source tree. At the earliest available source-code baseline,
2018-04-03, 193 of 1301 registered ordinary nodes were marked as deprecated, or
14.83\% (Table~\ref{tab:selected-snapshots}). By the KNIME Analytics Platform
5.0.0 source-date anchor on 2023-02-22, the count had increased to 433 of 1442
ordinary nodes, or 30.03\%. The current local source-date snapshot on June 28,
2026, contains 502 deprecated ordinary nodes out of 1506, or 33.33\%. These
figures show that the deprecated-node share approximately doubled between the
2018 baseline and the 2026 local snapshot.

The annual checkpoint table gives the same trend at regular intervals
(Table~\ref{tab:annual-snapshots}). The number of repositories with checkout
commits rises from 44 in 2018 to 90 in 2026, so early and late totals should
not be compared as if the mined source corpus were fixed. Nevertheless, the
deprecated-node share rises from 14.83\% in 2018 to 33.47\% at the 2026 annual
checkpoint. Hidden ordinary nodes remain much less frequent than deprecated
ordinary nodes, but they increase from 2 in 2018 to 24 in 2026. Deprecated
dynamic nodesets appear from 2020 onward and remain at 2 in the annual
checkpoints through 2026.

\begin{table}
\centering
\caption{Selected KNIME repository-mining snapshots.}
\label{tab:selected-snapshots}
\scriptsize
\begin{tabular}{lp{0.34\textwidth}rrr}
\toprule
Date & Version context & Nodes & Deprecated nodes &
\shortstack{Deprecated\\share} \\
\midrule
2018-04-03 & KNIME Analytics Platform 3.5.3 source-date anchor & 1301 &
193 & 14.83\% \\
2019-12-05 & KNIME Analytics Platform 4.1.0 source-date anchor & 1191 &
227 & 19.06\% \\
2023-02-22 & KNIME Analytics Platform 5.0.0 source-date anchor & 1442 &
433 & 30.03\% \\
2026-03-03 & KNIME Analytics Platform 5.11.0 source-date anchor & 1503 &
503 & 33.47\% \\
2026-06-28 & Current local source-date snapshot & 1506 & 502 & 33.33\% \\
\bottomrule
\end{tabular}
\end{table}

\begin{table}
\centering
\caption{Annual KNIME repository-mining checkpoint statistics.}
\label{tab:annual-snapshots}
\scriptsize
\begin{tabular}{llrrrrrr}
\toprule
Year & Snapshot date & Repositories & Nodes &
\shortstack{Deprecated\\nodes} &
\shortstack{Deprecated\\nodesets} &
\shortstack{Hidden\\nodes} &
\shortstack{Deprecated\\share} \\
\midrule
2018 & 2018-04-03 & 44 & 1301 & 193 & 0 & 2 & 14.83\% \\
2019 & 2019-01-01 & 45 & 1500 & 248 & 0 & 2 & 16.53\% \\
2020 & 2020-01-01 & 54 & 1200 & 228 & 2 & 2 & 19.00\% \\
2021 & 2021-01-01 & 60 & 1366 & 398 & 2 & 5 & 29.14\% \\
2022 & 2022-01-01 & 67 & 1424 & 386 & 2 & 8 & 27.11\% \\
2023 & 2023-01-01 & 72 & 1428 & 433 & 2 & 10 & 30.32\% \\
2024 & 2024-01-01 & 80 & 1434 & 436 & 2 & 9 & 30.40\% \\
2025 & 2025-01-01 & 86 & 1488 & 446 & 2 & 23 & 29.97\% \\
2026 & 2026-01-01 & 90 & 1509 & 505 & 2 & 24 & 33.47\% \\
\bottomrule
\end{tabular}
\end{table}

\begin{table}
\centering
\caption{Annual KNIME repository-mining transition statistics.}
\label{tab:annual-transitions}
\scriptsize
\begin{tabular}{lrrrrr}
\toprule
Year & Added & Removed &
\shortstack{Newly\\deprecated} &
\shortstack{No longer\\deprecated} &
\shortstack{Category\\changed} \\
\midrule
2018 & -- & -- & -- & -- & -- \\
2019 & 194 & 22 & 16 & 0 & 6 \\
2020 & 11 & 0 & 0 & 0 & 4 \\
2021 & 140 & 1 & 80 & 0 & 111 \\
2022 & 69 & 6 & 17 & 29 & 17 \\
2023 & 56 & 4 & 61 & 0 & 67 \\
2024 & 17 & 7 & 4 & 0 & 3 \\
2025 & 36 & 0 & 9 & 0 & 1 \\
2026 & 52 & 2 & 60 & 0 & 3 \\
\bottomrule
\end{tabular}
\end{table}

The transition table indicates that node deprecation changes occur in bursts
rather than at a constant rate (Table~\ref{tab:annual-transitions}). The largest
annual increases in newly deprecated node identities occur at the 2021, 2023,
and 2026 checkpoints, with 80, 61, and 60 newly deprecated node keys,
respectively. The 2022 checkpoint is the only annual checkpoint in which a
substantial number of common node keys are no longer marked deprecated: 29.
Category-path changes are also concentrated in a few years, especially 2021
and 2023. Because these transition counts use metadata-level node identity keys,
they should be read as evidence of source metadata evolution, not as direct
evidence that published workflows fail to open or execute.

These repository-level results directly address the second research question:
how KNIME node deprecation evolved over time. They also provide necessary
context for the workflow-level questions, because a published workflow that
depends on older node identities may encounter deprecated, hidden, removed, or
migrated components in a later KNIME version. However, the tables do not answer
how often uploaded or executed workflows use deprecated nodes. That question
requires workflow-level corpora and execution traces beyond repository
metadata.

Taken together, the evidence gives the intended picture of the study. First,
KNIME is visible in a broad and sustained scientific literature. Second,
influential papers that use or mention KNIME rarely provide the full
workflow-preservation package needed for later reproduction: executable
workflow files, KNIME versions, extension context, data, and code. Third, the
workflow experiment shows the practical consequence of this gap: 12 article
records yielded workflow artifacts or workflow directories, workflows from all
12 opened in the current local KNIME environment, and four article records had
at least one workflow execute successfully. Fourth, KNIME's source metadata
continues to evolve, and roughly one third of ordinary registered nodes are
marked deprecated in recent source snapshots. The result is not a claim that
all KNIME workflows fail. The result is a reproducibility risk: as KNIME
evolves, papers that preserve only workflow images or textual descriptions
become harder to diagnose, migrate, and reproduce.

\section{Discussion}

To our knowledge, this is the first focused empirical study connecting KNIME
workflow reporting in scientific publications with KNIME source-code evolution.
The contribution is intentionally narrower than a full reproducibility
benchmark. We do not estimate the failure rate of all published KNIME
workflows. Instead, we show a reproducibility risk pattern: KNIME is used in a
visible scientific literature, influential papers often do not publish
executable workflow artifacts, successful current-version execution was
confirmed for four of twelve downloaded article records in the manual
experiment, and the KNIME node ecosystem has accumulated a substantial
deprecated-node surface.

The results suggest three practical implications. First, a published workflow
image is not an executable artifact. Several highly cited papers provide
figures or textual descriptions of connected KNIME nodes, but this is not
enough to recover node settings, extension versions, input paths, or migration
requirements. Second, reporting a KNIME version is useful but incomplete.
Workflows may also depend on third-party extension update sites and node
libraries, as shown by the PAINS case where the RDKit workflow could be
updated and executed while the Indigo workflow was blocked by a missing
extension. Third, repository-level deprecation is a useful warning signal. It
does not prove that a workflow fails, but it identifies a compatibility surface
that grows as the platform evolves. Without careful workflow preservation and
version reporting, this growth makes old workflows harder to inspect and
repair.

The next development of this study should move from paper-level evidence to
workflow-level evidence. One direction is to create a larger corpus of public
KNIME workflows from supplements, myExperiment, KNIME Hub, GitHub, and other
repositories, then test whether they open in current KNIME versions and which
deprecated, hidden, migrated, or missing nodes they contain. A second direction
is to link deprecated nodes to migration metadata and replacement nodes in the
KNIME source tree, so that workflow warnings can distinguish low-risk
deprecations from nodes that require unavailable extensions or manual repair.
A third direction is to preserve reproducibility environments for selected
case studies by recording operating system, KNIME version, installed
extensions, update sites, input data, import warnings, execution errors, and
output differences.

Tools and services outside the present evidence base can provide the next
layer of usage evidence. \texttt{knime2py} can help identify KNIME workflows that users
try to translate or operationalize outside KNIME \cite{KaplanKnime2py2026},
and \texttt{k2pweb.org} can provide anonymized aggregate evidence about real uploaded
or executed workflows \cite{K2PWeb2026}. These
sources would address a question that the current paper deliberately leaves
open: how often researchers and analysts encounter deprecated KNIME nodes in
practice, not only how often deprecated nodes appear in the source repository
or in selected published artifacts.

\bibliographystyle{splncs04}
\bibliography{references}

\end{document}
